\documentclass[10pt,a4paper]{amsart}
\usepackage[margin=19mm]{geometry}
\usepackage{amsmath,amssymb,mathtools}
\usepackage[scr=rsfs,cal=euler]{mathalfa}
\usepackage{xfrac}
\usepackage[unicode,hidelinks,bookmarks=false]{hyperref}
\newcommand{\ie}{i.e.\ }
\newcommand{\cf}{cf.\ }
\newcommand{\resp}{respectively\ }
\newcommand{\Subsection}{\subsection}
\providecommand{\nobreakdash}{\nobreak\hbox{-}}
\setlength{\emergencystretch}{2em}
\multlinegap0pt
\usepackage{amssymb}
\usepackage{float}
\usepackage{mathtools}
\usepackage{tikz}
\usepackage{xcolor}
\newtheorem{lemma}{Lemma}
\newcommand{\PG}{\mathrm{PG}}
\title{The chromatic number of finite projective spaces}
\author{Anurag Bishnoi, Wouter Cames van Batenburg, Ananthakrishnan Ravi}
\date{Source: arXiv:2512.01760v3 (CC BY 4.0)}
\begin{document}
\maketitle

\noindent\emph{Source and licence.} The chromatic number of finite projective spaces. Version arXiv:2512.01760v3 (CC BY 4.0), distributed by arXiv under the Creative Commons Attribution 4.0 International licence, http://creativecommons.org/licenses/by/4.0/, as recorded in the arXiv licence metadata for this exact version and shown on the abstract page https://arxiv.org/abs/2512.01760v3. Full licence notice: \texttt{licenses/arxiv-2512.01760-CC-BY-4.0.txt}.\par\smallskip

\noindent\textit{Editorial scope.} Counted unit: the article's lemma lem:recursion (source0.tex lines 426--432). The article's complete original proof of it is delivered with it (source0.tex lines 434--535, one \texttt{proof} environment of 102 lines). No mathematics is written, repaired or re-derived: the statement and the proof are the article's own lines, in the article's own notation, and no retained line is substituted or re-flowed.  No further source-local context is needed: every cross-reference of every retained line resolves inside the two delivered spans.  Nothing was omitted from any retained span, so the package carries no omission record.  No retained line contains a citation, so the wrapper carries no bibliography and no bibliography entry.  No retained line contains a figure environment, an image-inclusion command or drawing code, so no image is delivered and the package declares no asset dependency.  Editorial formatting only, in addition: an AMS \texttt{amsart} preamble that restates the article's own declarations and macros used by the retained lines, namely the environments \texttt{lemma} on separate counters, together with the standard packages those lines need.  The retained environments print in this document as Lemma 1 in order of appearance, whereas the article prints the same items with the numbering of its own full document.  The complete native source of the article as distributed in the arXiv e-print for this exact version is bundled unchanged as \texttt{sources/190149/source0.tex}.
\begin{lemma}\label{lem:recursion}

For every prime power $q$ and every integer $1\le d\le n-1$,
\[
\chi_q(n)\le \chi_q(d)+ \chi_q(n+1-d)-1.
\]
\end{lemma}

\begin{proof}
Let $n-1\ge d\ge 1$.
Represent the points of $\PG(n-1,q)$ as $[v_1:\dots:v_n]$.
Define
\[
A \;=\; \bigl\{[v_1:\dots:v_n] : v_{n-d+1}=\dots=v_n=0\bigr\},
\]
\[
U \;=\; \bigl\{[v_1:\dots:v_n] : v_1=\dots=v_{n-d}=0\bigr\}.
\]
Then $A\cong\PG(n-d-1,q)$ and $U\cong\PG(d-1,q)$.
Every nonzero vector $v\in \mathbb F_q^n$ can be written uniquely as $v=(a,u)$, where $a\in \mathbb F_q^{n-d}$ and $u\in \mathbb F_q^d$. We write the corresponding projective point as $[a:u]$.
If $a$ is nonzero, then $[a:0]\in A$, and if $u$ is nonzero, then $[0:u]\in U$.

Take a proper $\chi_q(d)$-coloring $c_U$ of $U\cong\PG(d-1,q)$ and fix a reserved color $r$ of $c_U$.
Let $V_U$ be a set of vectors constructed from $U$ by fixing a particular coordinate vector for each point in $U$, say by taking the vector whose first non-zero coordinate is $1$.
Then for each vector $u \in \mathbb{F}_q^{d} \setminus \{0\}$, there exists a unique scalar $\lambda$ such that $u = \lambda v$ for some $v \in V_U$.
Define a function
\[
t:\mathbb{F}_q^{d} \longrightarrow \mathbb{F}_q
\]
by
\[
t(u) \;=\;
\begin{cases}
0, & u=0,\\
0, & u\neq 0 \text{ and } c_U([u])\neq r.\\
\lambda, & u = \lambda v  \text{ for some } v \in V_U \text{ and } c_U([u])=r,
\end{cases}
\]
By construction, for every $\lambda\in\mathbb F_q$ and every $u\in\mathbb F_q^d$, we have $t(\lambda u)=\lambda t(u)$.
Indeed, this is clear if $u=0$ or if $c_U([u])\neq r$. If $c_U([u])=r$ and $u=\mu v$ with $v\in V_U$, then $\lambda u=(\lambda\mu)v$, so $t(\lambda u)=\lambda\mu=\lambda t(u)$.

Now extend the underlying vector space of $A$, $\mathbb F_q^{n-d}$ to $\mathbb F_q^{n-d+1}$, and let $A^+\cong\PG(n-d,q)$ be the corresponding projective space.
Take a proper $\chi_q(n+1-d)$-coloring $c_{A^+}$ of $A^+$, using a palette disjoint from the colors of $c_U$.
Define
\[
c([a:u]) \;=\;
\begin{cases}
c_U([u]), & \text{if } u\neq 0 \text{ and } c_U([u])\neq r,\\[3pt]
c_{A^+}\bigl([a:t(u)]\bigr), & \text{otherwise}
\end{cases}
\]
Thus, $r$ is never used, and the total number of colors is at most $\chi_q(d)-1+\chi_q(n+1-d)$.
Moreover, this function is well-defined since $t(\lambda u) = \lambda t(u)$ which ensures that $[a:t(u)]$ is independent of the choice of representative for $[a:u]$.


We show that $c$ is a proper coloring.
Assume a line $\ell\subset\PG(n-1,q)$ is monochromatic in color $\alpha$.
Let $V_\ell\le \mathbb F_q^n$ be the $2$-dimensional vector subspace corresponding to the projective line $\ell$.
Define the linear projection $\varphi:\mathbb F_q^n=\mathbb F_q^{n-d}\oplus \mathbb F_q^d\longrightarrow \mathbb F_q^d$
by $\varphi(a,u)=u.$

\emph{Case 1: $\alpha$ is a color of $c_U$.}
By the definition of $c$, every point $[a:u]\in \ell$ satisfies $u\neq 0$ and $c_U([u])=\alpha\neq r.$
In particular, no nonzero vector of $V_\ell$ has $U$-component equal to $0$. Equivalently, $\ker(\varphi|_{V_\ell})=\{0\}.$
Thus $\varphi|_{V_\ell}$ is injective and its image is a $2$-dimensional vector subspace, which corresponds to a line of $U$. 
Every point of this line must have color $\alpha$ under $c_U$.
This gives a monochromatic projective line in $U$ for the coloring $c_U$, contradicting the fact that $c_U$ is proper.

\emph{Case 2: $\alpha$ is a color of $c_{A^+}$.}
Then for every $x=[a:u]\in\ell$ we have $u=0$ or $c_U([u])=r$, 
and 
$c(x)=c_{A^+}([a:t(u)])=\alpha$.
Again consider the projection map $\varphi$. 
If its image is a line of $U$, then it must be an $r$-colored line of $U$, contradicting the fact that $c_U$ is a proper coloring. 
If the projection is just the zero vector, then $t(u) = 0$ and we get a monochromatic line in $A^+$, which is again a contradiction. 
Therefore, the projection must be a single point of $U$, that is, a $1$-dimensional vector subspace.
Define $T: V_\ell \rightarrow \mathbb{F}_q^{n - d + 1}$ by $T(a, u) = (a, t(u))$.
We show that this map is a linear injective map and hence its image is also a $2$-dimensional vector space, thus giving us a monochromatic line of $A^+$, a contradiction.

Let $(a_1,u_1)$ and $(a_2,u_2)$ be two vectors in $V_\ell$. 
Since $\varphi(V_\ell)$ is a single $1$-dimensional vector subspace, there exists some $w\in V_U$ such that every $u\in \varphi(V_\ell)$ can be written uniquely as $u=sw$ for some $s\in\mathbb F_q$.
Thus $u_1=s_1w$ and $u_2=s_2w$ for some $s_1,s_2\in\mathbb F_q$.

Since every nonzero $u\in \varphi(V_\ell)$ corresponds to a point colored by the second rule, the point $[w]$ has color $r$ under $c_U$. Therefore, by the definition of $t$, if $u_i=s_iw$, then $t(u_i)=s_i$.
Say $T(a_1,u_1)=T(a_2,u_2)$. Then $a_1=a_2$ and $t(u_1)=t(u_2)$.
Since $t(u_1)=s_1$ and $t(u_2)=s_2$, this implies $s_1=s_2$, and hence $u_1=u_2$. Therefore $(a_1,u_1)=(a_2,u_2)$, showing that $T$ is injective.

We now show linearity.
Let $\lambda_1,\lambda_2\in\mathbb F_q$. Since $u_1=s_1w$ and $u_2=s_2w$, we have $\lambda_1u_1+\lambda_2u_2=(\lambda_1s_1+\lambda_2s_2)w$. Hence, by the definition of $t$, we get $t(\lambda_1u_1+\lambda_2u_2)=\lambda_1s_1+\lambda_2s_2=\lambda_1t(u_1)+\lambda_2t(u_2)$.
Therefore,
\[
T(\lambda_1(a_1,u_1)+\lambda_2(a_2,u_2))
=
(\lambda_1a_1+\lambda_2a_2,t(\lambda_1u_1+\lambda_2u_2))
=
\lambda_1T(a_1,u_1)+\lambda_2T(a_2,u_2).
\]
Thus $T$ is linear.

Since $T$ is linear and injective, $T(V_\ell)$ is a $2$-dimensional vector subspace of $\mathbb F_q^{n-d+1}$. Hence $T(V_\ell)$ corresponds to a projective line in $A^+$.
For every nonzero $(a,u)\in V_\ell$, the point $[a:u]$ lies on $\ell$, and by assumption it has color $\alpha$. Therefore $c_{A^+}([a:t(u)])=\alpha$, so every point of this line has color $\alpha$ under $c_{A^+}$.
This contradicts the properness of $c_{A^+}$.

\vspace{2ex}
\noindent
Both \textit{Case 1} and \textit{Case 2} lead to a contradiction, implying that $c$ must be a proper coloring of $\mathrm{PG}(n - 1, q)$, and thus
\[
\chi_q(n)\;\le\; (\chi_q(d)-1)+\chi_q(n+1-d)\;=\;\chi_q(d)+\chi_q(n+1-d)-1. \qedhere
\]
\end{proof}

\end{document}
